\section{Commercial systems, applications, and economics}
\label{sec:commercial}

\subsection{What the available evidence establishes}
The commercial landscape comprises distinct businesses: commissioned molten-glass objects, research printers using locally heated glass filament, and composite materials converted into glass after shaping. Their addressable parts, skill requirements, and qualification burdens differ. Consequently, a vendor list is not a ranking of interchangeable machines. This assessment separates supplier statements, experimentally characterized specimens, patent disclosures, and financial records. Public information was checked on 26 September 2026. No independently audited machine prices, production yields, or application-specific technology readiness levels were established; the qualitative maturity descriptions in Table~\ref{tab:commercial} should not be converted into numerical TRLs.

\begin{table}[htbp]
\centering
\caption{Commercial and industrial landscape. NR means not reported or not verified in the accessed primary sources. Operating temperatures are not maximum equipment ratings.}
\label{tab:commercial}
\scriptsize
\begin{tabularx}{\textwidth}{@{}P{0.14\textwidth}P{0.13\textwidth}P{0.17\textwidth}P{0.15\textwidth}Y@{}}
\toprule
Company / machine & Route & Maximum extrusion temperature & Feedstock & Target market; evidenced maturity \\
\midrule
Evenline / Mark III & Molten extrusion & NR; advertised operation $2000\,^{\circ}\mathrm{F}\approx1093\,^{\circ}\mathrm{C}$ & Studio, container, float glass & Art/design; operating studio platform. Structural masonry remains experimental \cite{com_evenline,com_masonry}. \\
Maple Glass Printing / Maple 4 in 2026 study & Heated-nozzle deposition & NR & Soda-lime rods in study & Research, education, architecture, art; experimental use verified, current sales unverified \cite{com_maple,acad_chhadeh2026}. \\
Nobula / Zero in 2024 study & Laser filament deposition & Current rated maximum NR & Fused-silica filament in study & Micro-optics; characterized research specimens and commercialization programme \cite{com_liu,com_nobulagrant}. \\
Glassomer / materials and parts service & Composite shaping and sintering & Not molten extrusion; room-temperature shaping & Silica--organic composite & Optics, microfluidics, prototypes; supplier materials and component offerings \cite{com_glassomer}. \\
UpNano + Glassomer / NanoOne + UpQuartz & Two-photon polymerization and sintering & Not molten extrusion; $1300\,^{\circ}\mathrm{C}$ sintering & Silica nanocomposite & Precision small parts; vendor process and resin offering \cite{com_upnano}. \\
\bottomrule
\end{tabularx}
\end{table}

\subsection{Company profiles and limits of comparison}
Evenline identifies Michael Stern as its founder and lead \cite{com_evenlinepeople}. Its technology page names Mark III, gives a $325\times325\times380$ mm fabrication envelope, and describes ceramic-nozzle deposition, manual separation of the object from the melt stream, and overnight cooling in a separate chamber. These details make operator involvement and annealer capacity relevant to throughput; the advertised envelope is not a guarantee for every geometry. A vendor process video is available at \href{https://www.evenline.co/technology}{Evenline's technology page}. Neither an equipment price nor a validated maximum extrusion temperature appears in that source \cite{com_evenline}.

The 2024 masonry paper characterizes G3DP3, an Evenline-operated system, and interlocking glass units \cite{com_masonry}. Its academic publication strengthens the evidence beyond photographs, but company-affiliated participation means this is not independent customer acceptance testing. The university account reports that the strongest units combined printed bodies with separately produced interlocking features, and presents larger self-supporting structures as future work \cite{com_mitbricks}. A tested prototype is therefore evidence of architectural potential, not an established building product. Qualification would need to address joints, impact, sustained and lateral loads, statistical variability, and installation details for the intended structure.

Maple remains identifiable through Monash University's venture directory and patent publication AU2021289614A1. The latter lists Nick Birbilis, Darren Robert Feenstra, and Tony Koutsonikolas as inventors and Maple Glass Printing Pty Ltd as applicant/assignee. Its disclosed system heats glass in a nozzle and describes driven solid feedstock; this establishes an invention disclosure, not throughput or machine reliability. Inventorship is not proof of a founding-team role, and the displayed patent-status field is not a legal opinion \cite{com_maple,com_maplepatent}. A 2026 study explicitly identifies Maple 4 and characterizes toolpath-dependent mechanical behavior of printed soda-lime specimens, providing recent experimental evidence \cite{acad_chhadeh2026}. The manufacturer's current website was not successfully verified, so present sales availability, price, and rated temperature ceiling remain unknown.

KTH's programme directory names Michael Fokine, Chunxin Liu, and Taras Oriekhov as Nobula's team. KTH Holding's May 2022 announcement confirms an investment and reports two patent filings, without establishing their subsequent grant status or disclosing investment value \cite{com_nobulateam,com_nobulainvest}. CORDIS verifies an EU contribution of EUR~2,499,875 for project 190122644, with total project cost EUR~3,571,250 and an end date of 30 April 2026. These are project-finance figures, not printer prices or cumulative company funding. Its product-performance language belongs to the project description and requires separate validation \cite{com_nobulagrant}.

Liu and colleagues used Nobula Zero to make fused-silica microlens arrays, with 200~$\mu$m filament and a CO$_2$ laser. Their study demonstrates a particular optical manufacturing workflow, including laser postprocessing; it cannot establish that every glass part requires no finishing. Several authors have Nobula affiliations. Supplier descriptions and this research should therefore be labelled separately, without calling either independent validation of a complete product range \cite{com_liu}. The live 2026 website confirms the company's laser-based manufacturing direction but did not expose a numerical specification table in the accessible page; current sales configuration, acceptance criteria, and price remain procurement questions \cite{com_nobula}.

Glassomer's history records its formation in Karlsruhe in 2018 and subsequent move to Freiburg. Its current offering includes composite materials, prototyping, and fused-silica components; this is a materials/process business rather than a like-for-like molten extrusion machine \cite{com_glassomerhistory,com_glassomer}. The UpNano collaboration extends the route to NanoOne two-photon polymerization using UpQuartz. Its 2024 manufacturer announcement specifies binder removal at $600\,^{\circ}\mathrm{C}$, sintering at $1300\,^{\circ}\mathrm{C}$, and approximately 30\% isotropic shrinkage. Those values apply to that reported material/process, not all Glassomer formulations. Published demonstrations include ferrules and a small filter, supporting precision-part targeting while leaving production yield unreported \cite{com_upnano}.

\subsection{Application fit and competitive alternatives}
The strongest near-term case for large molten beads is custom shape and visible material expression. Art, lighting, and limited-run packaging concepts can accept a layered surface when it is part of the design. Functional packaging adds closure tolerances, cleanability, impact resistance, and repeatability; an attractive vessel alone does not establish suitability for food, pharmaceuticals, or pressurization. Structural masonry is supported by specimen and assembly research, while broad architectural deployment remains an inference requiring application testing \cite{com_masonry}.

For lenses, optical channels, and chemical microdevices, high value per part can justify slower fabrication and metrology. Laser filament deposition has direct micro-optical evidence, and composite routes have supplier offerings in optics and microfluidics \cite{com_liu,com_glassomer}. Nevertheless, hermetic chemical vessels require verified leak tightness, pressure capability, connection integrity, and chemical compatibility for the actual fluid and temperature. Radioactive-waste vitrification containers are treated here as a proposed research direction: this source set establishes neither a qualified printed containment product nor a commercial deployment. Glass chemistry and a printable shape are insufficient evidence of long-term containment performance.

Conventional alternatives remain the baseline for purchasing decisions. Blowing and lampworking offer skilled manipulation of hollow forms; casting spreads tooling and furnace costs over repeated parts; machining and polishing support controlled surfaces on suitable blanks. Indirect AM introduces binder removal, shrinkage compensation, and sintering, whereas molten deposition transfers difficulty into thermal control, bead geometry, and annealing. These process differences motivate a matched-part comparison: identical composition or function, dimensional tolerances, inspected surface area, accepted yield, and annual volume. Comparing a millimetre optical component with a kilogram decorative print on cost per kilogram would conceal the product value and the inspection burden.

\subsection{Transparent cost model and illustrative sensitivities}
No supplier quotation is fabricated here. A reproducible planning model expresses cost per accepted part as
\begin{equation}
 C_{\mathrm{good}}=\frac{K/Y+F}{N_a y}
 +\frac{m c_m+E c_e+t_l c_l+c_q+c_p}{y},
\label{eq:commercialcost}
\end{equation}
where $K$ is installed capital, $Y$ depreciation life in years, $F$ annual fixed cost, $N_a$ annual attempted parts, and $y$ accepted yield. The remaining terms represent charged feedstock, electricity, labour, inspection, and consumables/postprocessing per attempt. All costs must share a currency and price basis. This simplified straight-line model excludes financing, taxes, residual value, and schedule-dependent rework; late failures are assumed to consume a full attempt. Parallel printer and annealer schedules must be checked before assuming $N_a$ is achievable.

For an explicitly hypothetical studio cell, let $K=\$100{,}000$, $Y=5$, $F=\$10{,}000$/year, $N_a=1000$, and $y=0.80$. Assume 2 kg of charged feedstock at \$2/kg, 20 kWh at \$0.15/kWh, one labour hour at \$30/h, \$10 inspection, and \$5 consumables per attempt. Annual fixed cost is \$30,000 and variable cost \$52/attempt, giving \$102.50/accepted part. These analyst-selected values are neither a price estimate for Evenline nor measured energy consumption. At 250 attempts/year, the same arithmetic gives \$215.00; at 2000 it gives \$83.75. At the baseline volume, lowering yield to 0.60 raises cost to \$136.67. Doubling electricity price raises baseline cost only to \$106.25 under these assumptions.

A second hypothetical small-part indirect-AM cell assumes $K=\$60{,}000$, $Y=5$, $F=\$6000$/year, $N_a=2000$, and $y=0.90$. Per attempt, allocate \$5 resin, 2 kWh at \$0.15/kWh, 0.25 labour hours at \$30/h, \$3 inspection, and \$2 firing consumables. The result is \$29.78/accepted part. This is not a finding that indirect AM is cheaper: the illustrative product is smaller, furnace energy is a batch allocation, and geometric acceptance may differ. Underfilled firing batches or additional optical polishing can reverse a superficial comparison. A useful quotation request therefore specifies the actual part, allowable defects, batch size, annual demand, training, furnace installation, maintenance, spare nozzles or optics, and acceptance tests. Until such matched quotations and meter readings exist, utilization and yield sensitivities are more defensible than a claimed market-wide CapEx or OpEx ranking.
